\SourceLeaf{269}%
den Abstand vom mittleren Breitenkreis und enthält außerdem die Abplattung oder das Quadrat von \(e\) als Faktor.

Dieser Vorteil lässt sich noch erheblich vergrößern, wenn man für die beliebige Funktion eine etwas abgeänderte, für die Rechnung fast ebenso bequeme Wahl trifft. Dazu zieht man eine zweite Konstante \(\alpha\) heran und setzt
\[
f\upsilon=\alpha\upsilon-i\log k
\]
Durch geeignete Wahl der beiden Konstanten kann man dann erreichen, dass die Abweichung des Vergrößerungsverhältnisses \(m\) vom Wert 1 in Bezug auf den Abstand vom mittleren Breitenkreis eine Größe dritter Ordnung ist, abgesehen vom Faktor \(ee\), der auch hier bestehen bleibt.

\Article{3}
\indent Bei dieser Wahl der Funktion \(f\) ergibt Formel 1 die folgenden Beziehungen.
\[
T=\alpha t \quad .\ .\ .\ .\ .\ .\ .\ .\ .\ .\ .\ .\ .\ .\ .\ .\ .\quad(2)
\]
\[
\operatorname{tang}\tfrac{1}{2}U=k\operatorname{tang}\tfrac{1}{2}w^\alpha(\tfrac{1+e\cos w}{1-e\cos w})^{\frac{1}{2}\alpha e} \quad .\ .\ .\ .\ .\ .\ .\ .\ .\quad(3)
\]
Für \(m\) erhält man aus den in der mehrfach erwähnten Schrift entwickelten Grundsätzen leicht den Ausdruck
\[
m=\frac{\alpha A\sin U\BareRoot{(1-ee\cos w^2)}}{a\sin w} \quad .\ .\ .\ .\ .\ .\ .\ .\ .\ .\quad(4)
\]
wobei \(a\) die große Halbachse des Ellipsoids und \(A\) den Radius der Kugel bezeichnet.

Differenziert man Gleichung 3 in logarithmischer Form, so ergibt sich
\[
\frac{\mathrm{d}U}{\sin U}=\frac{\alpha\mathrm{d}w}{\sin w}-\frac{\alpha ee\sin w\mathrm{d}w}{1-ee\cos w^2}
\]
oder
\[
\frac{\mathrm{d}U}{\mathrm{d}w}=\frac{\alpha(1-ee)\sin U}{(1-ee\cos w^2)\sin w} \quad .\ .\ .\ .\ .\ .\ .\ .\ .\ .\ .\ .\quad(5)
\]
Ebenso ergibt die Differentiation von Gleichung 4
\[
\begin{aligned}
\mathrm{d}\log m&=\operatorname{cotg}U\mathrm{d}U-\operatorname{cotg}w\mathrm{d}w+\frac{ee\cos w\mathord{.}\sin w\mathrm{d}w}{1-ee\cos w^2}\\
&=\operatorname{cotg}U\mathrm{d}U-\frac{(1-ee)\cos w\mathrm{d}w}{(1-ee\cos w^2)\sin w}
\end{aligned}
\]
